\chapter{Interprétabilité et gouvernance des modèles}\label{ch:ml:interpretability-and-model-governance}

Un comité des risques demande pourquoi le modèle a réduit de moitié sa position sur une action mardi. La réponse du chercheur, selon laquelle la
sortie du modèle à \omterm{def:ml:trees-and-boosting:boosting}{gradient boosting} a baissé, est vraie et inutile. Le comité veut savoir quelles variables ont bougé, si la
réponse du modèle à ces variables a du sens, s'il aurait dit la même chose le mois dernier, et qui a vérifié. Ce chapitre porte sur
la façon de répondre honnêtement à ces questions~: des explications justes sur le modèle même lorsque les variables sont corrélées,
des explications qui résistent à un réentraînement, et les documents (une \omterm{def:ml:interpretability-and-model-governance:card}{fiche de modèle}, un rapport de validation) par lesquels un modèle appris
entre dans la gestion du risque de modèle de l'établissement (livre~6, chapitre~26). Deux mesures l'ancrent. Avec deux variables corrélées à
0.99, la \omterm{def:ml:interpretability-and-model-governance:pd}{dépendance partielle} fausse l'effet d'une variable avec une erreur deux fois plus grande que celle des \omterm{def:ml:interpretability-and-model-governance:pd}{effets locaux accumulés}. Et sur un problème
doté d'un rapport signal sur bruit réaliste, les cinq variables les plus importantes selon la \omterm{def:ml:feature-engineering-selection-and-importance:shapley}{valeur SHAP} ne sont jamais les mêmes sur dix réentraînements.

\section{À quoi sert une explication}

\begin{definition}[Interprétabilité]\label{def:ml:interpretability-and-model-governance:interp}
L'\emph{interprétabilité}\index{interprétabilité} est le degré auquel une personne peut comprendre pourquoi un modèle produit ses
sorties~: globalement (comment la prédiction dépend de chaque variable sur l'ensemble des données) ou localement (pourquoi cette prédiction, pour cette
observation). Une explication est elle-même un modèle du modèle, et possède une précision propre.
\end{definition}

Les explications s'adressent à des publics différents. Un chercheur les utilise pour déboguer~: une variable qui compte et ne devrait pas compter est une fuite
(chapitre~3). Un gestionnaire des risques les utilise pour vérifier que les réponses du modèle ont le bon signe et la bonne ampleur, et que celui-ci ne
fasse rien d'absurde dans un marché qu'il n'a pas observé. Un comité ou un régulateur les utilise pour décider si l'établissement comprend
ce qu'il fait tourner. Rudin (2019) soutient que, pour les décisions à forts enjeux, un modèle interprétable doit être préféré à une boîte noire
assortie d'explications~; en trading, où les modèles sont souvent petits et le signal faible, c'est une option réelle (les \omterm{def:ml:linear-and-regularised-baselines:baseline}{modèles de référence} du chapitre~4),
et un modèle appris doit justifier sa complexité face à elle. Le modèle du chapitre est une régression par \omterm{def:ml:trees-and-boosting:boosting}{gradient boosting} (chapitre~5)
ajustée à $y = x_1 + x_2^2 + 2x_3x_4 + \varepsilon$ sur 6\,000 observations, où $x_2$ est, dans deux variantes, une copie bruitée de
$x_1$ (corrélations 0.92 et 0.99).

\section{Explications globales~: dépendance partielle et effets locaux accumulés}

\begin{definition}[Dépendance partielle, espérance conditionnelle individuelle, effets locaux accumulés]\label{def:ml:interpretability-and-model-governance:pd}
La \emph{dépendance partielle}\index{dépendance partielle} d'un modèle $f$ sur la variable $j$ est $\mathrm{PD}_j(v) = n^{-1}\sum_i
f(v, x_{i,-j})$~: la prédiction moyenne lorsque la variable $j$ est fixée à $v$ pour chaque observation (Friedman, 2001). Les courbes
d'\emph{espérance conditionnelle individuelle}\index{espérance conditionnelle individuelle} (ICE) sont les termes de cette
moyenne, une courbe par observation (Goldstein et coauteurs, 2015).
Les \emph{effets locaux accumulés}\index{effets locaux accumulés} (ALE) moyennent, dans des intervalles étroits de la variable $j$, la variation de la prédiction lorsque $j$ parcourt l'intervalle
pour les observations qui s'y trouvent, et cumulent ces variations (Apley et Zhu, 2020).
\end{definition}

La \omterm{def:ml:interpretability-and-model-governance:pd}{dépendance partielle} demande ce que le modèle prédit pour des combinaisons de variables qui peuvent ne jamais se présenter~: avec $x_2$ presque égal
à $x_1$, fixer $x_1 = -0.9$ pour une observation avec $x_2 = 0.9$ revient à interroger le modèle sur un point où il n'a aucune donnée, et un
ensemble d'arbres répond par ce que dit là son dernier découpage. Les ALE ne déplacent chaque observation que sur un court intervalle autour
de sa propre valeur (le \cref{lst:ml:gov:ale}). Le \Cref{tab:ml:gov:effects} mesure les deux par rapport à l'effet réel de $x_1$, une droite
de pente 1, et la \cref{fig:ml:gov:pdale} les montre à une corrélation de 0.99.

\begin{table}[htbp]
\centering\small
\begin{tabular}{@{}lccc@{}}
\toprule
& \multicolumn{3}{c}{corrélation de $x_1$ et $x_2$}\\
\cmidrule(l){2-4}
erreur quadratique moyenne par rapport à l'effet réel de $x_1$ & 0 & 0.92 & 0.99\\
\midrule
\omterm{def:ml:interpretability-and-model-governance:pd}{dépendance partielle} & 0.027 & 0.134 & 0.229\\
\omterm{def:ml:interpretability-and-model-governance:pd}{effets locaux accumulés} & 0.028 & 0.044 & 0.110\\
\bottomrule
\end{tabular}
\caption{Erreur de deux explications globales du même modèle ajusté, évaluée aux valeurs de $x_1$ présentes dans les données (l'effet réel
varie sur un intervalle de 2). Données~: \texttt{ml\_explain.effects}.}\label{tab:ml:gov:effects}
\end{table}

\begin{omfigure}
\begin{tikzpicture}
\begin{axis}[width=10.5cm,height=6cm,xlabel={$x_1$},ylabel={effet centré sur la prédiction},tick label style={font=\small},
  label style={font=\small},xmin=-1,xmax=1,grid=major,grid style={gray!25},legend cell align=left,
  legend style={font=\footnotesize,at={(0.02,0.98)},anchor=north west,fill=white,draw=none},table/col sep=comma]
\addplot[black,thick,dashed] table[x=x,y=truth]{figdata/ml/21-interpretability-and-model-governance/pdale.csv};
\addplot[omThm,very thick,mark=*,mark size=1.3pt] table[x=x,y=pd]{figdata/ml/21-interpretability-and-model-governance/pdale.csv};
\addplot[omDef,very thick,mark=square*,mark size=1.3pt] table[x=x,y=ale]{figdata/ml/21-interpretability-and-model-governance/ale.csv};
\legend{effet réel,dépendance partielle,effets locaux accumulés}
\end{axis}
\end{tikzpicture}

\omcaption{Effet de $x_1$ sur un modèle à \omterm{def:ml:trees-and-boosting:boosting}{gradient boosting} lorsque $x_2$ est une copie de $x_1$ de corrélation 0.99. Données~:
\texttt{ml\_explain.curves}.}\label{fig:ml:gov:pdale}
\end{omfigure}

Avec des variables indépendantes, les deux méthodes concordent. À mesure que la corrélation augmente, la \omterm{def:ml:interpretability-and-model-governance:pd}{dépendance partielle} absorbe une partie de l'effet quadratique de $x_2$
dans $x_1$, et son erreur est multipliée par cinq, puis par huit~; celle des ALE augmente aussi, car un modèle entraîné sur des variables presque
identiques ne peut pas dire lui-même quelle variable porte quel effet, mais elle reste à la moitié de celle de la \omterm{def:ml:interpretability-and-model-governance:pd}{dépendance partielle}. Aucune explication ne peut
séparer des effets que les données ne séparent pas~; l'ALE, au moins, n'ajoute pas d'erreurs issues de zones que les données n'ont jamais visitées.

Les courbes ICE montrent ce que cache une moyenne. La \omterm{def:ml:interpretability-and-model-governance:pd}{dépendance partielle} de $x_3$ est plate (pente 0.009)~: son effet passe entièrement par son
interaction avec $x_4$. Les courbes ICE de 30 observations (la \cref{fig:ml:gov:ice}) s'écartent en éventail, avec des pentes dont l'écart-type
vaut 0.91 et dont la corrélation avec le $x_4$ de l'observation est de 0.98~: le modèle a appris l'interaction, et la
moyenne de ses réponses l'a effacée.

\begin{omfigure}
\begin{tikzpicture}
\begin{axis}[width=10.5cm,height=6cm,xlabel={$x_3$},ylabel={prédiction},tick label style={font=\small},
  label style={font=\small},xmin=-0.9,xmax=0.9,grid=major,grid style={gray!25},table/col sep=comma]
\pgfplotsinvokeforeach{0,1,2,3,4,5,6,7,8,9,10,11,12,13,14,15,16,17,18,19,20,21,22,23,24,25,26,27,28,29}{
  \addplot[omProp!45,thin] table[x=x,y=c#1]{figdata/ml/21-interpretability-and-model-governance/ice.csv};
}
\addplot[black,ultra thick] table[x=x,y=pd]{figdata/ml/21-interpretability-and-model-governance/ice.csv};
\end{axis}
\end{tikzpicture}

\omcaption{Courbes ICE de $x_3$ pour 30 observations (fines) et leur moyenne (épaisse), pour le modèle aux variables indépendantes.
Données~: \texttt{ml\_explain.ice\_x3}.}\label{fig:ml:gov:ice}
\end{omfigure}

\begin{definition}[Substitut global]\label{def:ml:interpretability-and-model-governance:surrogate}
Un \emph{substitut global}\index{substitut global} (global surrogate) est un modèle interprétable (un arbre peu profond, un modèle linéaire) ajusté aux
prédictions d'un modèle complexe~; sa \emph{fidélité}, la part de la variance des prédictions qu'il reproduit, indique quelle part du
modèle qu'elle explique.
\end{definition}

Un arbre de profondeur trois ajusté aux prédictions du modèle ne découpe que selon $x_1$ et $x_2$ et reproduit 48~\% de leur variance~: il
explique la moitié additive du modèle et aucune part de l'interaction. Un substitut n'est une explication honnête qu'à condition
d'afficher sa fidélité juste à côté de lui.

\section{Explications locales et stabilité}

\begin{definition}[Explication contrefactuelle]\label{def:ml:interpretability-and-model-governance:cf}
Une \emph{explication contrefactuelle}\index{explication contrefactuelle} d'une prédiction est le plus petit changement des entrées qui
aurait modifié la décision, par exemple la valeur d'une variable pour laquelle la prévision change de signe (Wachter,
Mittelstadt et Russell, 2017).
\end{definition}

Les explications locales répondent à la question du comité au sujet du mardi. Les \omterm{def:ml:feature-engineering-selection-and-importance:shapley}{valeurs SHAP} (chapitre~6) répartissent une prédiction entre les
variables~; un contrefactuel dit ce qui aurait dû être différent. Pour une observation du modèle du chapitre, dont la
prédiction vaut 0.82 avec $x_1 = 0.46$, la prévision ne devient négative que si $x_1$ descend à $-0.44$~: l'essentiel de cette prédiction
provient de l'interaction, et $x_1$ seul devrait subir un changement important pour l'inverser.

Une explication qui change à chaque réentraînement explique la session d'entraînement, et non le marché. La \cref{fig:ml:gov:ice} et le
tableau ci-dessus proviennent d'un même ajustement~; le chapitre réentraîne aussi un modèle dix fois sur des rééchantillonnages bootstrap d'un problème comportant trois
variables fortes (coefficients 1, 0.7, 0.5), cinq variables faibles (0.2 chacune) et douze sans effet, puis classe les variables selon la
valeur absolue moyenne des SHAP à chaque fois (\cref{lst:ml:gov:stab}). Lorsque les variables expliquent 66~\% de la variance, les cinq
premières variables forment le même ensemble dans 22~\% des réentraînements, et le top~5 du premier réentraînement conserve son ordre avec une
corrélation de Kendall moyenne de 0.89~: les trois variables fortes sont toujours en tête, et celles des deux variables faibles qui complètent la liste relèvent d'un tirage au sort
parmi cinq ex æquo. Lorsqu'elles expliquent 5~\%, plus proche d'une prévision de rendement, les cinq premières ne forment jamais le même ensemble et la
corrélation de Kendall tombe à 0.78~; dans certains réentraînements, une variable nulle entre dans le top~5. Un rapport qui cite «~les cinq variables les plus
importantes du modèle~» à partir d'un seul ajustement rapporte du bruit sur ses deux dernières places.

\section{Documentation~: la fiche de modèle}

\begin{definition}[Fiche de modèle]\label{def:ml:interpretability-and-model-governance:card}
Une \emph{fiche de modèle}\index{fiche de modèle} (model card) est un document court qui accompagne un modèle et précise son objectif, ses données
d'entraînement, son évaluation, sa performance par tranches pertinentes, son comportement, ses limites et ses approbations (Mitchell et coauteurs, 2019)~; dans une société
elle constitue l'entrée du modèle appris dans l'inventaire des modèles (livre~6, chapitre~26).
\end{definition}

La fiche est générée, non rédigée~: ses champs sont renseignés à partir de la session d'entraînement et du rapport de validation, et le générateur
refuse de déclarer une fiche complète tant que des champs manquent. La fiche du modèle de démonstration du chapitre compte quinze champs et
signale deux champs manquants, le suivi et les approbations, ce qui correspond exactement à l'état d'un modèle construit et validé mais non
déployé. Son champ «~limites~» consigne le seul fait qu'un utilisateur doit connaître~: lorsque $x_1$ et $x_2$ sont corrélées, il faut lire l'ALE, et non la
\omterm{def:ml:interpretability-and-model-governance:pd}{dépendance partielle}.

\begin{dated}{2026-09}{Règles relatives au risque de modèle et modèles appris}\label{dat:ml:interpretability-and-model-governance:rules}
Les orientations américaines sur le risque de modèle ont été révisées le 17 avril 2026 (SR 26-2, remplaçant SR 11-7)~; le livre~6 (chapitre~26)
les décrit toutes deux. La déclaration de supervision SS1/23 de la Prudential Regulation Authority, \emph{Model risk management principles for
banks}, publiée le 17 mai 2023, fixe cinq principes et les applique à tous les modèles utilisés pour éclairer des décisions d'entreprise,
«~quelle que soit la technologie~», y compris «~l'usage de l'intelligence artificielle dans les techniques de modélisation telles que l'apprentissage automatique
dans la mesure où il s'applique à l'usage des modèles en général~»~; sa version en vigueur est entrée en application le 23 avril 2026,
après les modifications de faible portée de la PRA d'avril 2026.
\end{dated}

\section{Validation d'un modèle appris}

La validation d'un modèle appris pose les questions du livre~6 (cohérence conceptuelle, analyse des résultats, comparaison à un indice de référence,
remise en cause effective) ainsi que quelques questions que les modèles appris rendent plus pressantes~: la fuite d'information, la stabilité d'un réentraînement à l'autre, et la
parité entre le modèle validé et celui qui tourne en production. La \cref{tab:ml:gov:checklist} est la liste de contrôle du chapitre pour son modèle de démonstration, renseignée là où les preuves existent et marquée
«~non exécuté~» là où elles manquent~; les points numériques passent par le
\texttt{firm.modelval} du livre~6.

\begin{table}[htbp]
\centering\small
\begin{tabular}{@{}p{4.3cm}lp{7.0cm}@{}}
\toprule
item & résultat & preuve\\
\midrule
performance hors échantillon & réussi & $R^2$ 0.860 sur 2\,000 nouvelles observations\\
comparaison au challenger & réussi & régression ridge $R^2$ 0.349 sur les mêmes observations\\
parité d'implémentation & réussi & les arbres exportés (chapitre~5) reproduisent le modèle à l'identique sur 200 lignes\\
analyse des résultats & réussi & 27 résultats sur 250 hors de l'intervalle à 90~\% (25 attendus), queue binomiale 0.37\\
cohérence conceptuelle, fuite d'information, stabilité, explications, suivi & non exécuté & à renseigner par le validateur\\
\bottomrule
\end{tabular}
\caption{La liste de contrôle de validation du modèle du chapitre. Données~: \texttt{ml\_explain.validation}.}\label{tab:ml:gov:checklist}
\end{table}

\begin{method}[Expliquer et gouverner un modèle appris]\label{met:ml:gov:method}
\begin{enumerate}
\item Expliquer les effets par l'ALE lorsque les variables sont corrélées, afficher les courbes ICE lorsque des interactions sont soupçonnées, et indiquer la
fidélité d'un substitut avec le substitut.
\item Présenter les classements de variables issus de plusieurs réentraînements, avec leur stabilité~; ne nommer que les variables stables.
\item Répondre aux questions sur des décisions isolées par les \omterm{def:ml:feature-engineering-selection-and-importance:shapley}{valeurs SHAP} et les contrefactuels calculés sur le modèle déployé.
\item Générer la \omterm{def:ml:interpretability-and-model-governance:card}{fiche de modèle} à partir des enregistrements d'entraînement et de validation~; l'enregistrer dans l'inventaire des modèles~; laisser le validateur
renseigner la liste de contrôle, et ne pas déployer avec des points non exécutés.
\end{enumerate}
\end{method}

\section{Tutoriel~: expliquer le modèle au comité des risques}\label{tut:ml:interpretability-and-model-governance:tut}

\begin{tutorial}
\textbf{Objectif.} Expliquer un modèle boosté avec PD, ICE, ALE et un substitut~; mesurer la stabilité de ses classements SHAP~; générer
sa fiche et sa liste de contrôle. \textbf{État final~:} \cref{tab:ml:gov:effects,tab:ml:gov:checklist},
\cref{fig:ml:gov:pdale,fig:ml:gov:ice}.
\begin{tutsteps}
\item \textbf{\omterm{def:ml:interpretability-and-model-governance:pd}{Effets locaux accumulés}.}
\omcode{code/firm/modelcard/firm_modelcard.py}{50}{69}{ALE du premier ordre avec des classes en quantiles.}\label{lst:ml:gov:ale}
\item \textbf{Stabilité des classements d'un réentraînement à l'autre.}
\omcode{code/firm/modelcard/firm_modelcard.py}{103}{112}{Accord sur les ensembles de tête et corrélation de Kendall entre réentraînements.}\label{lst:ml:gov:stab}
\item \textbf{Exécuter} \texttt{ml\_explain.effects()}, \texttt{ice\_x3()}, \texttt{surrogate()}, \texttt{stability(1.0)},
\texttt{stability(6.0)}, \texttt{validation()}, \texttt{card().render()} et \texttt{fig\_explain.py}.
\end{tutsteps}
\textbf{Pour aller plus loin.} Ajouter l'ALE du second ordre pour la paire $(x_3, x_4)$~; remplacer les réentraînements bootstrap par
des réentraînements sur des années successives~; ajouter à la liste de contrôle un canari de fuite d'information (chapitre~3).
\end{tutorial}

\section{Construction~: explications et fiches de modèle}\label{bld:ml:interpretability-and-model-governance:modelcard}

\begin{build}
\bfield{Rôle} Des explications dont la précision et la stabilité sont mesurées, et un modèle appris documenté et validé comme
tout autre modèle.
\bfield{Interface} \texttt{partial\_dependence}, \texttt{ice}, \texttt{ale}, \texttt{global\_surrogate},
\texttt{counterfactual}, \texttt{rank\_stability}~; \texttt{ModelCard} avec \texttt{render} et \texttt{missing}~;
\texttt{benchmark\_check}, \texttt{outcome\_check}, \texttt{validation\_checklist} (sur \texttt{firm.modelval}).
\bfield{Règles} Aucune explication sans sa précision ou sa fidélité~; aucun classement sans sa stabilité~; aucune fiche avec des champs manquants
déclarée complète~; les points de liste de contrôle sans preuve sont «~non exécutés~».
\bfield{Tests d'acceptation} \texttt{code/firm/modelcard/tests/}~: PD et ALE retrouvent une fonction additive connue~; l'ALE est juste et
la PD fausse sur une paire corrélée avec un piège hors variété~; les pentes ICE retrouvent une interaction~; un substitut d'un arbre est parfait~;
les contrefactuels franchissent le seuil~; la stabilité des classements vaut 1 pour des classements identiques~; la fiche liste ses champs manquants.
\bfield{Pour aller plus loin} ALE du second ordre~; valeurs d'interaction SHAP~; fiches dans l'outil de suivi d'expériences (chapitre~25).
\end{build}

\begin{omsources}
\item J.~H.~Friedman, ``Greedy function approximation: a gradient boosting machine'', \emph{Annals of Statistics} 29(5), 2001.
\item A.~Goldstein, A.~Kapelner, J.~Bleich et E.~Pitkin, ``Peeking inside the black box: visualizing statistical learning with
plots of individual conditional expectation'', \emph{Journal of Computational and Graphical Statistics} 24(1), 2015.
\item D.~W.~Apley et J.~Zhu, ``Visualizing the effects of predictor variables in black box supervised learning models'',
\emph{Journal of the Royal Statistical Society B} 82(4), 2020.
\item M.~Mitchell and co-authors, ``Model cards for model reporting'', \emph{FAT*}, 2019.
\item S.~Wachter, B.~Mittelstadt et C.~Russell, ``Counterfactual explanations without opening the black box: automated
decisions and the GDPR'', SSRN 3063289, 2017.
\item C.~Rudin, ``Stop explaining black box machine learning models for high stakes decisions and use interpretable models
instead'', \emph{Nature Machine Intelligence} 1, 2019.
\item Prudential Regulation Authority, SS1/23, \emph{Model risk management principles for banks}, 2023 (version en vigueur
2026).
\end{omsources}

\section{Exercices}

\begin{exercise}[$\star$]\label{exo:ml:interpretability-and-model-governance:1}
Pour $f(x_1, x_2) = x_1 + x_2^2$ avec $x_1$ et $x_2$ indépendantes et uniformes sur $[-1, 1]$, écrire la \omterm{def:ml:interpretability-and-model-governance:pd}{dépendance partielle} en
$x_1$. Qu'est-ce qui change si $x_2 = x_1$ exactement et que le modèle est $f$ lui-même~?
\end{exercise}

\begin{exercise}[$\star$]\label{exo:ml:interpretability-and-model-governance:2}
Avec 250 résultats et un intervalle à 90~\%, combien d'exceptions attend-on~? 27 est-il alarmant~?
\end{exercise}

\begin{exercise}[$\star$]\label{exo:ml:interpretability-and-model-governance:3}
Un \omterm{def:ml:interpretability-and-model-governance:surrogate}{substitut global} a une fidélité de 0.48. Que peut-on dire et que ne peut-on pas dire à partir de cette valeur~?
\end{exercise}

\begin{exercise}[$\star\star$]\label{exo:ml:interpretability-and-model-governance:4}
Pourquoi la \omterm{def:ml:interpretability-and-model-governance:pd}{dépendance partielle} de $x_3$ est-elle plate alors que $x_3$ compte~?
\end{exercise}

\begin{exercise}[$\star\star$]\label{exo:ml:interpretability-and-model-governance:5}
Pourquoi les quatrième et cinquième places du classement SHAP changent-elles d'un réentraînement à l'autre, même avec un signal fort~?
\end{exercise}

\begin{exercise}[$\star\star$]\label{exo:ml:interpretability-and-model-governance:6}
\emph{Trouver la faille.} «~Les \omterm{def:ml:feature-engineering-selection-and-importance:shapley}{valeurs SHAP} du modèle montrent que la variable de révision des bénéfices explique 40~\% de la prévision, donc le
modèle est un modèle de révisions et le comité peut le considérer comme tel.~»
\end{exercise}

\begin{exercise}[$\star\star\star$]\label{exo:ml:interpretability-and-model-governance:7}
\emph{Programmation.} Relancer l'expérience de stabilité au niveau de signal élevé avec dix réentraînements sur des dixièmes disjoints des données
(400 observations chacun) au lieu de rééchantillonnages bootstrap des 4\,000 observations. La stabilité augmente-t-elle ou diminue-t-elle, et pourquoi~?
\end{exercise}

\begin{exercise}[$\star\star\star$]\label{exo:ml:interpretability-and-model-governance:8}
Montrer que, pour un modèle additif $f = \sum_jf_j(x_j)$, l'ALE de la variable $j$ est égal à $f_j$ à une constante près, quelle que soit
la dépendance entre les variables, et que la \omterm{def:ml:interpretability-and-model-governance:pd}{dépendance partielle} l'est aussi. Où la \omterm{def:ml:interpretability-and-model-governance:pd}{dépendance partielle} se trompe-t-elle alors~?
\end{exercise}

\section{Problème~: expliquer le modèle au comité des risques}

\begin{problem}[{Problème du week-end --- une explication a une barre d'erreur}]\label{pb:ml:interpretability-and-model-governance:1}
Le modèle du chapitre, les réentraînements, la fiche et la liste de contrôle.

\medskip\noindent\textbf{Partie I --- Effets globaux.}
\begin{enumerate}
\item Quelles sont les erreurs de la PD et de l'ALE à chaque corrélation~?
\item Pourquoi l'erreur de l'ALE augmente-t-elle aussi~?
\item Que montrent les courbes ICE de $x_3$ que la PD ne montre pas~?
\item Qu'explique le substitut, et que manque-t-il~?
\end{enumerate}

\medskip\noindent\textbf{Partie II --- Local et stable.}
\begin{enumerate}[resume]
\item Que dit le contrefactuel de la prévision de l'exemple~?
\item Quelle est la stabilité du top cinq à chaque niveau de signal~?
\item Que doit contenir un rapport sur l'importance des variables~?
\item Comment répondre à la question du comité sur mardi~?
\end{enumerate}

\medskip\noindent\textbf{Partie III --- Gouvernance.}
\begin{enumerate}[resume]
\item Que contient la \omterm{def:ml:interpretability-and-model-governance:card}{fiche de modèle}, et que manque-t-il~?
\item Quels points de la liste de contrôle sont validés, et sur quelle preuve~?
\item Pourquoi la parité d'implémentation figure-t-elle sur la liste~?
\item Que dit l'encadré daté sur la façon dont les superviseurs traitent les modèles appris~?
\end{enumerate}

\medskip\noindent\textbf{Partie IV --- Le verdict.}
\begin{enumerate}[resume]
\item Énoncer le \emph{résultat nommé}~: l'erreur de la \omterm{def:ml:interpretability-and-model-governance:pd}{dépendance partielle} par rapport à l'effet réel sous variables corrélées, comparée à
celle de l'ALE, et la stabilité de rang du top cinq des variables SHAP sur dix réentraînements.
\item Faut-il remplacer ici le modèle boosté par un modèle interprétable~?
\item Dans quel niveau (livre~6) faudrait-il classer un modèle d'alpha appris, et pourquoi~?
\item À quoi ressemble une remise en cause effective pour un modèle appris~?
\item Que faut-il revalider après un réentraînement~?
\item Comment explique-t-on un modèle dont les entrées sont des plongements (chapitre~10)~?
\item Que ne doit jamais contenir une \omterm{def:ml:interpretability-and-model-governance:card}{fiche de modèle}~?
\item En une phrase~: qu'est-ce qu'une explication~?
\end{enumerate}
\end{problem}

\section{Questions d'entretien}

\begin{interviewq}[$\star$ \iqroles{chercheur, ingénieur ML}]\label{iq:ml:interpretability-and-model-governance:1}
Quelle est la différence entre la \omterm{def:ml:interpretability-and-model-governance:pd}{dépendance partielle} et l'ALE, et quand cela importe-t-il~?
\end{interviewq}

\begin{interviewq}[$\star\star$ \iqroles{chercheur}]\label{iq:ml:interpretability-and-model-governance:2}
Quelle est la stabilité des classements de variables fondés sur SHAP, et comment la mesureriez-vous~?
\end{interviewq}

\begin{interviewq}[$\star\star$ \iqroles{ingénieur ML}]\label{iq:ml:interpretability-and-model-governance:3}
Que contient une \omterm{def:ml:interpretability-and-model-governance:card}{fiche de modèle} pour un modèle de trading~?
\end{interviewq}

\begin{interviewq}[$\star\star$ \iqroles{chercheur, trader}]\label{iq:ml:interpretability-and-model-governance:4}
Un gérant de portefeuille demande pourquoi le modèle a vendu une action. Que montrez-vous~?
\end{interviewq}

\begin{interviewq}[$\star\star$ \iqroles{ingénieur ML, chercheur}]\label{iq:ml:interpretability-and-model-governance:5}
Comment valideriez-vous un modèle boosté avant la mise en production~?
\end{interviewq}

\begin{interviewq}[$\star\star\star$ \iqroles{chercheur}]\label{iq:ml:interpretability-and-model-governance:6}
Quand préféreriez-vous un modèle interprétable à une boîte noire expliquée en trading~?
\end{interviewq}
